% Clean standalone section for later integration into main.tex.
% This file deliberately uses the all-loose formulation of Kaloshin's mixed
% system.  It avoids any extra transversality assumption on rigid levels.

\subsection{An unconditional fixed-itinerary theorem}\label{subsec:unconditional-fixed-itinerary}

We now close the Khovanskii--Rolle argument for an arbitrary cyclic crossing
itinerary without imposing an auxiliary admissibility hypothesis.  Consider a
fixed transversal itinerary with $k\ge2$ flights.  On its $j$th flight assume
\begin{equation}\label{eq:uf-model}
 x'=a_jx+b_j(t),\qquad
 b_j(t)=\sum_{m=-M}^{M}\beta_{j,m}e^{imt},
 \qquad a_j\in\mathbb R,\quad M\ge1.
\end{equation}
Let $c_j,c_{j+1}$ be the two switching values at the endpoints of the flight.
Write $q_j=e^{it_j}$ and let $\tau_j\in(0,2\pi)$ be the positively oriented gap
from $q_j$ to $q_{j+1}$, with $q_{k+1}=q_1$ and
$\sum_j\tau_j=2\pi$.  Introduce one positive private variable per flight,
\begin{equation}\label{eq:uf-private}
 y_j=\begin{cases}
 R_j=e^{a_j\tau_j},&a_j\ne0,\\
 T_j=\tau_j,&a_j=0.
 \end{cases}
\end{equation}
For $a_j\ne0$ the endpoint matching equation is
\begin{equation}\label{eq:uf-Q-nonzero}
 Q_j=c_{j+1}-R_jc_j-
 \sum_{m=-M}^{M}\frac{\beta_{j,m}}{im-a_j}
 \bigl(q_{j+1}^m-R_jq_j^m\bigr)=0,
\end{equation}
while for $a_j=0$ it is
\begin{equation}\label{eq:uf-Q-zero}
 Q_j=c_{j+1}-c_j-\beta_{j,0}T_j-
 \sum_{m\ne0}\frac{\beta_{j,m}}{im}
 \bigl(q_{j+1}^m-q_j^m\bigr)=0.
\end{equation}
Thus every $Q_j$ is polynomial--Laurent and
\begin{equation}\label{eq:uf-local-support}
 \Newt(Q_j)\subseteq
 L_j:=[-M,M]_{q_j}+[-M,M]_{q_{j+1}}+[0,1]_{y_j}.
\end{equation}

Let $\mathcal C_k$ be one connected component of the configuration space of
$k$ labelled cyclically ordered distinct points on $S^1$, and put
\begin{equation}\label{eq:uf-M}
 \mathcal M_k=\mathcal C_k\times(0,\infty)^k.
\end{equation}
The physical exponential/time relations are the integral manifold of the
$k$ one-forms
\begin{equation}\label{eq:uf-eta}
 \eta_j=\begin{cases}
 d\log y_j-a_j\,d\tau_j,&a_j\ne0,\\
 dy_j-d\tau_j,&a_j=0.
 \end{cases}
\end{equation}

\begin{lemma}[Global separating chain and covering]\label{lem:uf-covering}
The physical flight locus is a separating solution of
$\eta_1=\cdots=\eta_k=0$ on $\mathcal M_k$.  Moreover
\begin{equation}\label{eq:uf-rho}
 \rho_k(q,y)=\prod_{j=1}^k
 \left(2-\frac{q_{j+1}}{q_j}-\frac{q_j}{q_{j+1}}\right)
 \frac{y_j}{(1+y_j)^2}
\end{equation}
is a covering function in the sense of Khovanskii--Rolle theory.  For every
regular $\varepsilon>0$, its level is represented, after clearing positive
denominators, by the Laurent polynomial
\begin{equation}\label{eq:uf-Peps}
 P_\varepsilon=
 \prod_{j=1}^k
 \left(2-\frac{q_{j+1}}{q_j}-\frac{q_j}{q_{j+1}}\right)y_j
 -\varepsilon\prod_{j=1}^k(1+y_j)^2.
\end{equation}
\end{lemma}

\begin{proof}
For $a_j\ne0$ set $\phi_j=\log y_j-a_j\tau_j$, and for $a_j=0$ set
$\phi_j=y_j-\tau_j$.  Then $\eta_j=d\phi_j$, and $d\phi_j$ has a nonzero
component in the private $y_j$ direction.  Consequently the chain
\[
 \Gamma_r=\{\phi_1=\cdots=\phi_r=0\},\qquad r=1,\ldots,k,
\]
is smooth, and $\Gamma_r$ is the cooriented boundary in $\Gamma_{r-1}$ of
$\{\phi_r<0\}\cap\Gamma_{r-1}$.  This is a separating chain.

The phase factor in \eqref{eq:uf-rho} is
\[
 2-e^{i\tau_j}-e^{-i\tau_j}=4\sin^2(\tau_j/2)>0.
\]
The closure of $\mathcal C_k$ in $(S^1)^k$ is compact.  Hence a sequence in
$\mathcal M_k$ with no accumulation point must either approach a collision
$\tau_j\to0$ for some $j$, or satisfy $y_j\to0$ or $y_j\to\infty$ for some
$j$.  In all cases \eqref{eq:uf-rho} tends to zero.  Therefore it is a covering
function.  The polynomial \eqref{eq:uf-Peps} is obtained by multiplying the
level equation by $\prod_j(1+y_j)^2$.
\end{proof}

For the support estimate define the common box
\begin{equation}\label{eq:uf-S}
 S=[-2M,2M]^k_q\times[0,2]^k_y\subset\mathbb R^{2k}.
\end{equation}
Then
\begin{equation}\label{eq:uf-base-containment}
 \sum_{j=1}^kL_j\subseteq S,
 \qquad
 \Newt(P_\varepsilon)\subseteq S.
\end{equation}
Indeed, each phase coordinate occurs in exactly two local flight boxes, while
in \eqref{eq:uf-Peps} its exponent lies between $-2$ and $2$; the latter range
is contained in $[-2M,2M]$ because $M\ge1$.  Each private exponent in the
covering polynomial lies between $0$ and $2$.

Use on $\mathcal M_k$ the mixed coframe consisting of $d\arg q_1,\ldots,
d\arg q_k$, $d\log y_j$ on nonzero-slope flights, and $dy_j$ on zero-slope
flights.  In this coframe the coefficients of every $\eta_j$ are constant.
Logarithmic differentiation in $q$ and in nonzero-slope $y$ preserves Laurent
exponents, while ordinary differentiation in a zero-slope $y_j=T_j$ only lowers
its nonnegative exponent.  Therefore all differential coefficient supports
remain inside the boxes in \eqref{eq:uf-base-containment}.

For an integer $N\ge1$ define
\begin{align}
 A_0(N)&=2^{N(N-1)/2},\label{eq:uf-A0}\\
 A_j(N)&=
 2^{\frac{(j-1)(j-2)}2+\frac{(N-j)(N-j-1)}2}
 (2^{j-1}+1)^{N-j},\qquad 1\le j\le N,\label{eq:uf-Aj}\\
 \Theta_N&=A_0(N)+\frac12\sum_{j=1}^{N}A_j(N).
 \label{eq:uf-Theta}
\end{align}

\begin{theorem}[Unconditional sparse toric fixed-itinerary bound]
\label{thm:uf-fixed-itinerary}
For the constant-slope finite-Fourier equation \eqref{eq:uf-model}, the number
$N_X(k,M)$ of isolated crossing limit cycles realizing a prescribed cyclic
$k$-flight itinerary satisfies
\begin{equation}\label{eq:uf-fixed-bound}
 \boxed{
 N_X(k,M)\le
 \Theta_{2k}\,(2k)!\,(8M)^k.}
\end{equation}
The estimate is uniform in the switching levels and in all coefficients
$a_j,\beta_{j,m}$, and includes zero slopes.
\end{theorem}

\begin{proof}
We first count nondegenerate solutions of the transition equations.  Regard the
$k$ matching equations
\[
 Q=(Q_1,\ldots,Q_k)=0
\]
as the loose equations in Kaloshin's mixed functional--Pfaffian system, with
no rigid equations.  This is precisely the $s=0$ case in the definition of the
Khovanskii number.  The separating solution is the physical flight locus from
\cref{lem:uf-covering}.

A nondegenerate transition solution is a regular point of the restriction of
$Q$ to this separating leaf.  Given any finite collection of such solutions,
the inverse function theorem gives pairwise disjoint neighborhoods in which all
sufficiently small target values $b\in\mathbb R^k$ have one continued solution.
By Sard's theorem choose such a $b$ to be a regular value of $Q$ on the
separating leaf.  Hence the cardinality of the chosen finite collection is
bounded by the Khovanskii number of the all-loose mixed system.  This
regularization shifts the $Q_j$ by constants only; because the origin belongs
to every $L_j$, no Newton box is enlarged.

Kaloshin's reduction replaces the $k$ loose equations by the exact one-forms
$dQ_1,\ldots,dQ_k$ and then eliminates the resulting $2k$ one-forms
\[
 \eta_1,\ldots,\eta_k,dQ_1,\ldots,dQ_k.
\]
The general Khovanskii--Rolle reduction gives one all-contact algebraic system
and, for each of the $N=2k$ reduction steps, one system in which the covering
function is inserted at that step, with coefficient $1/2$.

Consider first the all-contact branch.  If $H_r$ is the new contact equation at
step $r$, \eqref{eq:uf-base-containment} and the support-preserving mixed
coframe give inductively
\begin{equation}\label{eq:uf-contact-recursion}
 \Newt(H_r)\subseteq2^{r-1}S,
 \qquad r=1,\ldots,N.
\end{equation}
Indeed, the still uneliminated original $dQ_j$ together contribute at most
$S$, while each previously created rigid contact equation contributes its own
Newton box.  Thus the dilation factors obey
$a_1=1$ and $a_r=1+\sum_{\ell<r}a_\ell=2^{r-1}$, whose product is $A_0(N)$.

If the covering equation is inserted at step $j$, its Newton polytope is in
$S$.  Before that step the factors are $1,2,\ldots,2^{j-2}$; at the covering
step the factor is $1$; afterwards the same support recurrence gives
\[
 a_r=2^{r-j-1}(2^{j-1}+1),\qquad r>j.
\]
The product of all $N$ dilation factors is exactly $A_j(N)$ in
\eqref{eq:uf-Aj}.

Every final auxiliary system consists of $2k$ Laurent polynomials in
$(q_1,\ldots,q_k,y_1,\ldots,y_k)\in(\mathbb C^*)^{2k}$.  The box $S$ contains
the origin, so subtracting the small regular levels required during the
Khovanskii--Rolle induction does not enlarge any adopted Newton box.  Likewise,
clearing the positive denominators in the covering equation only multiplies its
differential by a positive factor on the covering level and does not change the
relevant real intersection set.  Bernstein's theorem, monotonicity and
multilinearity of mixed volume give
\[
 \operatorname{MV}(a_1S,\ldots,a_NS)
 =\left(\prod_{r=1}^Na_r\right)(2k)!\operatorname{Vol}_{\rm Eucl}(S).
\]
Since
\[
 \operatorname{Vol}_{\rm Eucl}(S)=(4M)^k2^k=(8M)^k,
\]
the all-contact system contributes at most
$A_0(N)(2k)!(8M)^k$ and the $j$th covering system contributes at most
$A_j(N)(2k)!(8M)^k$.  Applying the Khovanskii--Rolle coefficients proves
\eqref{eq:uf-fixed-bound} for nondegenerate transition solutions.

Finally, the scalar additive-bias perturbation used for crossing cycles by
Bravo and Trinidad-Forte replaces $x'=S(t,x)$ by $x'=S(t,x)+\lambda$ and
shows that, for an upper bound, isolated crossing cycles may be replaced by
simple crossing cycles whose transition solutions are nondegenerate.  The
perturbation changes only the zero Fourier coefficients and leaves all constant
slopes and all Newton boxes above unchanged.  Transversality preserves the
prescribed itinerary.  Therefore the same uniform estimate bounds all isolated
crossing limit cycles with that itinerary.
\end{proof}

\begin{corollary}[Constant-slope Hilbert-type bound]
\label{cor:uf-global}
Let $\mathcal H_{\rm cs}^{\times}(n,M)$ be the supremum of the number of
crossing limit cycles of scalar periodic piecewise-linear equations with $n+1$
regions, constant slopes, and trigonometric forcing of degree at most $M\ge1$.
Then
\begin{equation}\label{eq:uf-global-crossing}
 \mathcal H_{\rm cs}^{\times}(n,M)
 \le 2Mn
 \max_{2\le k\le2Mn}
 \bigl[\Theta_{2k}(2k)!(8M)^k\bigr].
\end{equation}
Following the global bookkeeping of Bravo and Trinidad-Forte, the corresponding
bound after including the at most two constant-sign outer cycles is the
right-hand side of \eqref{eq:uf-global-crossing} plus $2$.  If the piecewise
vector field is continuous, their tangency-removal argument shows that this
latter bound applies to all isolated limit cycles, including tangential ones.
\end{corollary}

\begin{proof}
Bravo and Trinidad-Forte prove that a crossing limit cycle has at most $2Mn$
crossings and that there are at most $2Mn$ crossing sequences up to cyclic
permutation.  Apply \cref{thm:uf-fixed-itinerary} to each sequence.  Their
global proof adds at most two constant-sign cycles outside the extreme
switching levels.  In the continuous case their additive-bias argument removes
tangencies without increasing the number of regions or the Fourier degree and,
in the present subclass, without changing the constant slopes.
\end{proof}

\begin{remark}[Asymptotic gain over the general Pfaffian encoding]
For fixed $k$, the new bound grows as $C_kM^k$.  The dominant
Khovanskii--Rolle dilation in \eqref{eq:uf-fixed-bound} is
\[
 A_0(2k)=2^{k(2k-1)},
\]
whereas the direct Pfaffian construction for time-dependent slopes uses a chain
of length $4k$ and hence carries the general Khovanskii factor
$2^{(4k)(4k-1)/2}=2^{2k(4k-1)}$.  Thus the constant-slope toric lift changes the
Fourier-degree dependence from the general $M^{4k}$ scale to $M^k$ and reduces
the leading quadratic exponent in $k$ from $8k^2$ to $2k^2$.  The sharper
interaction-graph mixed-volume theorem remains useful for further reductions
which retain some of the original flight equations as rigid equations.
\end{remark}
